\documentclass[10pt,twocolumn]{article}
\usepackage[margin=1.9cm,columnsep=0.7cm]{geometry}
\usepackage{lmodern}\usepackage{mathptmx}
\usepackage{graphicx,booktabs,amsmath,amssymb}
\usepackage[hang,flushmargin]{footmisc}
\usepackage{caption}\captionsetup{font=small,labelfont=bf}
\usepackage{titlesec}\titleformat{\section}{\large\bfseries}{\thesection}{0.6em}{}
\usepackage{fancyhdr}\pagestyle{fancy}\fancyhf{}
\fancyfoot[C]{\tiny Tese Parte 2 — Etrização Computacional · github.com/camillanapoles/quest003-prion-v127 · tiers [SIM]/[ORGANOID]/[MOUSE]/[HUMAN]}
\fancyfoot[R]{\thepage}\renewcommand{\headrulewidth}{0pt}
\usepackage{xcolor}\usepackage{microtype}

\title{\bfseries ETRIZAÇÃO COMPUTACIONAL EM DOENÇAS PRIÔNICAS: APLICADA À PLATAFORMA TERAPÊUTICA PrP-V127}
\author{\normalsize Camilla N.\thanks{Open Prion \& Molecular Engineering Consortium. AI sob supervisão humana contínua; não é autora (SW-S01/03). Gated: guardian.py --profile part2.}}
\date{}

\begin{document}\maketitle

\begin{abstract}\small
A Parte 1 demonstrou a tese de design (plataforma V127, $\theta^{*}{=}0{,}333$, reproduzida em dois ambientes). A Parte 2 --- este documento --- é a tese realizada e formaliza o método: \textbf{etrização computacional} (\emph{Computational Etrization}), passos P0--P6. \textbf{Eixo:} tudo aqui é simulação e é dito que é; não prometemos validar --- se dados reais forem análogos, os passos seguintes já estão dados (antecipação \emph{bancada}). \textbf{Base de Validade} mandatória: simulação não substitui o laboratório; linhagem completa por dado. \textbf{[SEM ANO]:} prognóstico de simulação opera em tempo futuro. Resultados-da-tese = achados [SIM] executados: $\theta^{*}$; regras; sensibilidade discriminadora; estimador calibrado; tabela-decisão.
\end{abstract}

\section*{Nota à Leitura: sobre o termo Etrização}
Este documento introduz \textbf{etrização} como neologismo científico composicional intencional. O radical \emph{etr-} evoca o \emph{être} (potencial de vir a ser) e o \emph{atrium} (antecâmara virtual$\to$real). \textbf{Etrização} nomeia o estado do conhecimento que esta tese produz: sistema operacionalizado computacionalmente a partir de dados reais publicados, estruturalmente robusto, \textbf{ainda não validado empiricamente}. Como o `átomo' nomeou antes da prova um objeto que fundou ciência, etrização nomeia o processo que produz objetos-conceito operacionais na saúde. O método é a própria etrização (P0--P6, Cap.3); a operação técnica é a parametrização com proveniência (\S3.2). Cinco critérios distinguem etrização de mera simulação --- cadeia de dados, transposição justificada, incerteza quantificada, escopo delimitado, transição condicionada --- \textbf{todos cumpridos nesta tese}.

\section{Introdução}
\textbf{Problema:} pesquisa translacional em doenças raras paralisa-se: sem dado clínico não há financiamento; sem previsão, o dado se desperdiça (seis fracassos priônicos). \textbf{Questões:} Q1 formalizar método com rigor experimental? Q2 produzir resultados próprios? Q3 posicionamento na família? \textbf{Objetivos:} OE1 nomear/formalizar etrização P0--P6; OE2 Base de Validade; OE3 entregar resultados [SIM]; OE4 continuidade como método; OE5 revisão hostil. \textbf{Hipóteses:} H1 forma experimental-análoga; H2 predições falseáveis travadas; H3 posição estrutural distinta.

\section{Fundamentação}
\textbf{Gargalo:} seis falhas clínicas registry-bound (quinacrina, doxiciclina, PPS, flupirtina, PRN100, minociclina). \textbf{Mecanismo} validado em 4 níveis (população$\to$camundongo$\to$culturas$\to$AAV). \textbf{In-silico trials} simulam o ensaio; etrização simula a continuação da pesquisa. \textbf{\S2.5 Fundamento epistemológico:} etrização como operação ontológica --- a hierarquia de evidência da saúde não nomeia o conhecimento potencial-verificável; etrização preenche essa lacuna.

\section{Metodologia}
\textbf{P0} identificação (SLR auditada; linhagem) $\to$ \textbf{P1} parametrização (proveniência por parâmetro) $\to$ \textbf{P2} execução (determinística, self-tests: massa 100\%, Thiele 0,5\%) $\to$ \textbf{P3} colheita (critérios pré-declarados) $\to$ \textbf{P4} prognóstico (travado por release) $\to$ \textbf{P5} pesquisa derivada $\to$ \textbf{P6} confronto opcional.
\textbf{Equações:} $\partial(\alpha c)/\partial t = \nabla \cdot (D_{\text{ef}} \nabla c) - \nabla \cdot (\mathbf{v}c) + S(\mathbf{x}) - k_{\text{ef}}c$; freeS$=(1+\kappa c)^{-2}$; $\theta \equiv (1+\kappa c_{\text{pico}})^{-1}$; duplicação $\approx$12,1\,d; 1 unidade $=$ 144\,d.
\textbf{Base de Validade:} tríade não-substitui/prognóstico-antecipação/rigor-convertido + linhagem dos 6 dados.
\textbf{Mapeamento análogo:} participantes$\to$fontes; instrumentos$\to$código; coleta$\to$runs; prontuário$\to$registro.

\section{Resultados}
$\theta^{*}{=}0{,}333$; regras 1--3 (anel 8--12\,mm; malha $\xi\ge5r_p$; redose $\le$7\,d); sensibilidade discriminadora (C$_{50}$ 10$\times$ insensível; forma funcional falseável por dose-resposta); estimador $\theta_{\text{obs}}$ calibrado (bias $-0{,}008$ na fronteira; v1.1 rejeitada); tabela-decisão (biomassa $>$ raio saturado).

\section{Achados, Impactos, Áreas Correlatas; Próximas Abordagens}
\textbf{Impactos:} direto (DCJ/E200K-BR); plataforma (AD/PD $>$50M); meta (etrização demonstrada); epistemológico (legitima forma autônoma de saber). \textbf{Correlatas:} farmacologia de raras; oncologia; prion-like; toxicologia/regulação. \textbf{Declaração mandatória:} base simulada; laboratório necessário; dados humanos independentes de valores $\Rightarrow$ parametrizam contra obtidos $=$ passo à frente.

\section{Discussão}
O que significam: limiar comparável; design quantitativo inédito; 3 frentes falseáveis. Não significam: contenção medida; eficácia clínica; substituição do laboratório.

\section{Conclusões por Objetivo}
\textbf{OE1--OE5 alcançados.} \textbf{H1--H3} corroboradas. \textbf{Fecho epistemológico:} ``Como o átomo tornou pesquisável o invisível, a etrização torna pesquisável o ainda-não-medido.''

\section*{Referências}
38 fontes ABNT (ver registro). Concordância 54 claims$\to$refs (tabela no documento mestre).

\end{document}
